\section{Alternative formalisms}
\label{sec:alternative-formalisms}

FRFP is one point in a broader design space of human--AI collaboration architectures. This section
positions FRFP against competing formalisms, clarifies which alternatives do \emph{not} reduce to
FRFP (because they step outside the FRFP axiom class), and explains why FRFP is the appropriate
choice for the theorem-level aims of this volume.

\subsection{The design space: what varies across architectures}

Any protocol-level formalism must decide (at minimum):
\begin{enumerate}[leftmargin=2em]
  \item \textbf{Where correctness lives.} Is correctness a property decidable in explicit space (a predicate or oracle on artefacts), or a property of tacit endorsement states?
  \item \textbf{Whether tacit is writable.} Is there a channel \(T\to E\) that exports acceptance states into machine-manipulable specifications, policies, or constraints?
  \item \textbf{Who may access tacit.} May AI systems condition on tacit state (or modify it), or is tacit strictly human-exclusive?
  \item \textbf{How acceptance is structured.} Is there a single acceptance boundary, or multiple gates (individual, committee, institution) with distinct authority and standards?
  \item \textbf{How uncertainty is represented.} Is endorsement discrete (a finite judgment quotient), ordered (preorders), or probabilistic (graded belief and risk thresholds)?
\end{enumerate}
FRFP fixes these choices in a specific way: correctness is localized in tacit space; AI is confined to
explicit moves; the explicit--tacit boundary is one-way \(E\to T\) with no return channel \(T\to E\);
and the tacit interface is intentionally minimal so that the explicit algebra carries the technical
structure.

\subsection{Architectures that do \emph{not} reduce to FRFP}

Because FRFP is proved (in Appendix~A) to be \emph{initial} only within a category of frameworks that
share the FRFP epistemic primitives and axioms, an alternative architecture avoids ``reducing to
FRFP'' precisely by violating at least one of those axioms. The following are genuine competitors in
that sense.

\paragraph{(A) Bidirectional boundary architectures (\(T\leftrightarrow E\)).}
These formalisms introduce an essential write-back channel \(F:T\to E\), by which tacit endorsement
states externalize into explicit policies, specifications, rubrics, or machine-checkable constraints.
This matches institutional learning patterns (postmortems \(\to\) checklists \(\to\) standards), but it
is outside FRFP by construction: FRFP forbids \(T\to E\). In such systems, the acceptance criteria are
not merely applied at review points; they are continuously rewritten into the explicit layer.

\paragraph{(B) Tacit-in-the-loop AI.}
Here AI systems may condition on tacit state or participate in tacit updates (e.g., morphisms of the
form \(T\times E\to E\) or AI-mediated operators \(T\to T\)). This architecture is attractive for
personalization and closed-loop assistance, but it breaks human-exclusive correctness and blurs
responsibility: the system can implicitly steer endorsement by reading or shaping tacit states. This
is therefore not an instance of FRFP.

\paragraph{(C) Explicit-only correctness architectures.}
In some regimes, correctness can be reduced to a complete explicit oracle: proof checking, model
checking, exhaustive tests under a closed-world specification. In such settings, ``acceptance'' is an
explicit computation on artefacts, not a tacit judgment. This is again outside FRFP's target class:
FRFP assumes a genuine tacit layer that is not exhausted by any fixed explicit oracle.

\paragraph{(D) Multi-boundary acceptance architectures.}
Real institutions often implement layered gates of endorsement (individual sign-off, peer review,
risk committee approval, regulatory clearance). A formalism may represent these as multiple tacit
layers and multiple boundaries, with cross-talk between gates. If the ontology class fixes a single
boundary primitive, then such multi-gate architectures do not reduce to FRFP without extending the
primitive structure.

\paragraph{(E) Probabilistic tacit semantics.}
A competing formalism may model tacit state as a richer belief object (e.g., probability measures,
utility-weighted acceptance thresholds, risk budgets) rather than a preorder with a finite observable
judgment quotient. When the tacit algebra differs materially (graded endorsement rather than discrete
quotients and idempotent operators), the architecture falls outside the specific tacit-interface axioms
used in FRFP's initiality claims.

\subsection{Which alternatives can be viewed as generalizations of the scientific method}

Several competing architectures can still be regarded as ``generalizations of the scientific method,''
but they generalize different facets of it.

\begin{itemize}[leftmargin=2em]
  \item \textbf{Explicit-only correctness} generalizes the formal-verification slice of science (mathematics, parts of theoretical computer science), where correctness is internal to a fixed calculus.
  \item \textbf{Bidirectional \(T\leftrightarrow E\)} generalizes the meta-scientific process by which tacit know-how becomes codified into explicit methods, standards, and tooling.
  \item \textbf{Multi-boundary acceptance} generalizes the institutional reality of science as layered endorsement (lab judgment, journal review, replication, community uptake).
  \item \textbf{Probabilistic tacit semantics} generalizes science as graded belief and risk management under uncertainty.
\end{itemize}
FRFP generalizes a particular core: \emph{stabilization by accountable review}, where public explicit
artefacts are strengthened by explicit checks and then cross a responsibility-bearing acceptance
boundary.

\subsection{Why FRFP among the alternatives}

FRFP is chosen to match a specific problem class and to enable theorem-level analysis.

\paragraph{Target regime.}
We aim at long-horizon workflows in which (i) requirements evolve, (ii) explicit artefacts and
evaluation criteria drift across extended traces, and (iii) no fixed explicit oracle exhausts what it
means for an artefact to be adequate. In this regime, architectures that place correctness entirely in
\(E\) are aiming at a different class of problems.

\paragraph{Accountability and non-circular semantics.}
FRFP enforces that endorsement is attributable: AI systems can produce and transform explicit
artefacts, but only humans occupy correctness-bearing tacit states. The one-way boundary \(E\to T\)
prevents implicit rewriting of acceptance criteria by the system itself. Architectures with essential
\(T\to E\) write-back or tacit-in-the-loop AI can be valuable, but they entangle acceptance and
generation in ways that complicate responsibility and undermine the invariance and non-automation
results that this volume targets.

\paragraph{Minimality and canonicity within the axiom class.}
Within the class of formalisms that share FRFP's epistemic primitives and constraints (explicit/tacit
separation, one-way boundary, human-exclusive correctness, AI confined to explicit moves), FRFP is
initial: any competing framework satisfying the same assumptions admits a unique structure-preserving
interpretation from FRFP. In this precise sense, FRFP is not an arbitrary design choice inside its
assumption class; alternatives that keep the same primitives are instances/refinements rather than
routes around the theorem-level consequences.

\subsection{Practical reading of this comparison}

The existence of non-FRFP architectures is not a defect: it clarifies scope. FRFP is a protocol-level
formalism for domains where correctness-bearing endorsement cannot be reduced to explicit computation,
and where the central risks are long-horizon drift, accountability loss, and explicit-only governance
limits. Competing architectures are appropriate when one wishes to (i) treat correctness as fully
explicit, (ii) explicitly codify tacit judgments into machine-checkable policy, (iii) close the loop
between system behavior and human internal state, or (iv) model layered institutional gates or graded
belief as first-class. Subsequent volumes may study carefully delimited departures from FRFP's axioms;
the present volume fixes them to make the foundational theorems sharp.
